% !TEX program = xelatex
\documentclass[aspectratio=169,11pt]{beamer}
\usepackage{fontspec}
\IfFontExistsTF{Latin Modern Roman}{\setmainfont{Latin Modern Roman}}{\setmainfont{Times New Roman}}
\usepackage{graphicx,booktabs,amsmath,amssymb,xcolor}
\usetheme{default}
\usecolortheme{default}
\usefonttheme{serif}
\usefonttheme{professionalfonts}
\definecolor{AcademicBlue}{RGB}{51,51,179}
\setbeamercolor{normal text}{fg=black,bg=white}
\setbeamercolor{structure}{fg=AcademicBlue}
\setbeamercolor{title}{fg=AcademicBlue}
\setbeamercolor{frametitle}{fg=AcademicBlue,bg=white}
\setbeamercolor{block title}{fg=AcademicBlue,bg=white}
\setbeamercolor{block body}{fg=black,bg=white}
\setbeamerfont{normal text}{family=\rmfamily,size=\normalsize}
\setbeamerfont{title}{family=\rmfamily,size=\LARGE,series=\mdseries}
\setbeamerfont{frametitle}{family=\rmfamily,size=\Large,series=\mdseries}
\setbeamerfont{author}{family=\rmfamily,size=\large}
\setbeamerfont{institute}{family=\rmfamily,size=\normalsize}
\setbeamerfont{date}{family=\rmfamily,size=\normalsize}
\setbeamerfont{block title}{family=\rmfamily,size=\normalsize,series=\mdseries}
\setbeamerfont{block body}{family=\rmfamily,size=\small}
\setbeamerfont{footline}{family=\rmfamily,size=\scriptsize}
\setbeamertemplate{navigation symbols}{}

\setbeamertemplate{itemize item}{%
    \small\raise0.15ex\hbox{$\bullet$}%
}

\setbeamertemplate{itemize subitem}{%
    \scriptsize\raise0.15ex\hbox{$\bullet$}%
}

\setbeamertemplate{footline}{%
    \hfill
    \usebeamerfont{footline}%
    \usebeamercolor[fg]{structure}%
    \insertframenumber\,/\,\inserttotalframenumber
    \hspace{1em}\vspace{0.6em}%
}

\setbeamersize{
    text margin left=0.7cm,
    text margin right=0.7cm
}
\renewcommand{\arraystretch}{1.16}
\title{Pharmacies at the Hospital Gate:\\Drug Markup Regulation and Spatial Concentration}
\author{Wenxia Chen}
\institute{Xiamen University}
\date{November 11, 2026}
\newcommand{\allresultsroot}{D:/Hospital_Pharmacy program/results/results_v6_slides_1111}
\newcommand{\paperresultsroot}{\allresultsroot}
% Change this path only if your output directory has moved.
\newcommand{\resultsroot}{../data/hospital_pharmacy/county_pharm_panel/results_county_v4}
\newcommand{\hospresultsroot}{../data/hospital_pharmacy/hosp_pharm_panel/results_v5_all_public}
\newcommand{\resultimage}[2]{%
  \IfFileExists{\resultsroot/#1_regression/#2}{%
    \includegraphics[width=\linewidth,height=0.58\textheight,keepaspectratio]{\resultsroot/#1_regression/#2}%
  }{%
    \fbox{\parbox[c][0.42\textheight][c]{0.90\linewidth}{%
      \centering\scriptsize Missing image:\\
      Sample: \texttt{#1}\\Check the image path in the source.}}%
  }%
}
\newcommand{\eventimage}[1]{%
  \IfFileExists{\resultsroot/#1}{%
    \includegraphics[height=0.67\textheight,keepaspectratio]{\resultsroot/#1}%
  }{%
    \fbox{\parbox[c][0.50\textheight][c]{0.85\linewidth}{%
      \centering\scriptsize Missing full-sample event-study image.}}%
  }%
}
\newcommand{\countywideimage}[1]{%
  \IfFileExists{\resultsroot/#1}{%
    \includegraphics[width=\linewidth,height=0.72\textheight,keepaspectratio]{\resultsroot/#1}%
  }{%
    \fbox{\parbox[c][0.55\textheight][c]{0.90\linewidth}{%
      \centering\scriptsize Missing county-level image.}}%
  }%
}
\newcommand{\countypairimage}[1]{%
  \IfFileExists{\resultsroot/#1}{%
    \includegraphics[width=\linewidth,height=0.66\textheight,keepaspectratio]{\resultsroot/#1}%
  }{%
    \fbox{\parbox[c][0.50\textheight][c]{0.90\linewidth}{%
      \centering\scriptsize Missing county-level image.}}%
  }%
}
\newcommand{\hospresultimage}[1]{%
  \IfFileExists{\hospresultsroot/#1}{%
    \includegraphics[width=\linewidth,height=0.54\textheight,keepaspectratio]{\hospresultsroot/#1}%
  }{%
    \fbox{\parbox[c][0.42\textheight][c]{0.90\linewidth}{%
      \centering\scriptsize Missing hospital-level image.}}%
  }%
}
\newcommand{\hosplargeimage}[1]{%
  \IfFileExists{\hospresultsroot/#1}{%
    \includegraphics[width=0.92\textwidth,height=0.72\textheight,keepaspectratio]{\hospresultsroot/#1}%
  }{%
    \fbox{\parbox[c][0.55\textheight][c]{0.85\textwidth}{%
      \centering\scriptsize Missing hospital-level image.}}%
  }%
}
\newcommand{\hospquadimage}[1]{%
  \IfFileExists{\hospresultsroot/#1}{%
    \includegraphics[width=\linewidth,height=0.335\textheight,keepaspectratio]{\hospresultsroot/#1}%
  }{%
    \fbox{\parbox[c][0.26\textheight][c]{0.90\linewidth}{%
      \centering\scriptsize Missing hospital-level image.}}%
  }%
}
\newcommand{\threepanels}[2]{%
  \begin{columns}[T,onlytextwidth]
    \begin{column}{0.325\textwidth}
      \centering\small Full sample (bz)\\[0.5em]\resultimage{bz}{#1}
    \end{column}
    \begin{column}{0.325\textwidth}
      \centering\small AOI sample\\[0.5em]\resultimage{AOI}{#1}
    \end{column}
    \begin{column}{0.325\textwidth}
      \centering\small Open sample\\[0.5em]\resultimage{open}{#1}
    \end{column}
  \end{columns}
  \par#2
}
\newcommand{\sectiondivider}[1]{%
  \begin{frame}
    \vfill
    \centering
    {\usebeamercolor[fg]{structure}\LARGE #1}
    \vfill
  \end{frame}%
}

\begin{document}
\begin{frame}[plain]
  \titlepage
\end{frame}
\sectiondivider{Background and Data}
\begin{frame}{Institutional Background}
China's pharmaceutical reform had two key components. These components
have potentially ambiguous implications for retail pharmacies.
\begin{itemize}\setlength\itemsep{1.2em}
  \item[$\uparrow$] \textbf{Separated prescribing from dispensing:} allowing
patients to fill prescriptions at retail pharmacies rather than purchasing drugs exclusively from
hospitals.
  \item[$\downarrow$] \textbf{Reduced drug prices in public hospitals:} by requiring them to sell medications
at procurement cost, eliminating the previous markup on drug sales.
\end{itemize}
\end{frame}
\begin{frame}{Reform Rollout and Scope}
\vspace{-0.95em}
\small
The reform began with county-level public-hospital pilots and expanded nationwide in stages.
\begin{itemize}
  \setlength\itemsep{1.2em}
  \setlength\topsep{1.2em}
  \item In 2012, the first pilot covered roughly 300 counties or county-level cities.
  \item By September 2017, public hospitals nationwide had removed drug mark-ups, except for designated traditional Chinese medicine products.
\end{itemize}
\end{frame}
\begin{frame}{Data}
\small
\begin{itemize}\setlength\itemsep{0.65em}
 \item Business-Registration: Hospitals records, ownership, opening and cancellation dates, and geographic coordinates. 
  \item DZDP: Pharmacies records and locations.
  \item Official government documents and related reports: county policy years.
  \item 2,512 counties and 110,547 hospital-year observations covering 2013--2018.
  \item Counties appear in every year, before the centralized drug procurement policy announced in November 2018 and piloted from January 2019.
\end{itemize}
\end{frame}

\sectiondivider{Framework and Empirical Strategy}

\begin{frame}{Conceptual Framework}
\small
A pharmacy at distance $d$ from hospital $h$ serves residential demand $A(d)$ and
hospital-related demand $F_hs(d)\theta$. With margin $m$ and fixed costs and rent $c(d)$,
\[
\pi_h(d)=m\bigl[A(d)+F_hs(d)\theta\bigr]-c(d). \tag{1}
\]
The reform raises the outside prescription share from $\theta$ to $\theta+\delta$ and
lowers the retail margin from $m$ to $m-\kappa$. Holding competition and rents fixed,
\[
\Delta\pi_h(d)=F_hs(d)\bigl[(m-\kappa)\delta-\kappa\theta\bigr]-\kappa A(d). \tag{2}
\]
\begin{itemize}\setlength\itemsep{0.45em}
 \item The prescription gain scales with patient volume $F_h$ and declines with distance through $s(d)$.
 \item The residential margin loss does not depend on the hospital.
 \item Pharmacy entry should shift toward the gates of large hospitals and away from surrounding neighborhoods.
\end{itemize}
\end{frame}

\begin{frame}{Hospital-Level Sample, Treatment, and Distance Rings}
\small
\begin{itemize}\setlength\itemsep{0.5em}
 \item 28,989 hospitals; 110,547 hospital-years, 2013--2018, in the balanced set of 2,512 counties.
 \item All public hospitals, including public TCM hospitals, are treated from the first calendar year after county adoption. Private hospitals are comparisons.
 \item Three disjoint bands: gate (0--200 m), near ring (200--500 m), outer ring (500--1,000 m). Also report the 0--1,000 m total.
\end{itemize}
\[
 NetEntry_{htr}=Stock_{htr}-Stock_{h,t-1,r}. \tag{3}
\]
\footnotesize Net entry is undefined in the first hospital-year and after a gap.
The preferred lag-valid regression sample has 74,924 observations; stock regressions have 103,766.
\end{frame}

\begin{frame}{Summary Statistics}
\centering\scriptsize\setlength{\tabcolsep}{5.5pt}
{\renewcommand{\arraystretch}{1.02}\input{table_summary_statistics_1028.tex}}
\par\vspace{0.35em}
\raggedright\tiny
Notes: The sample contains the hospital-years in the net-entry regressions. Pharmacy counts are active Dianping pharmacies within each radius, and net entry is the annual change in that count. Tier shares use hospitals with a classified tier.
\end{frame}

\begin{frame}{Baseline Model}
\small
\[
 Y_{hctr}=\beta_r\bigl(Public_h\times Post_{ct}\bigr)
             +\alpha_h+\lambda_{ct}+\varepsilon_{hctr}. \tag{4}
\]
\begin{itemize}\setlength\itemsep{0.65em}
 \item $Y_{hctr}$ is net pharmacy entry or the pharmacy stock in distance band $r$ around hospital $h$ in county $c$ and year $t$.
 \item $Post_{ct}=\mathbb{1}\{t>PolicyYear_c\}$.
 \item Hospital FE $\alpha_h$ absorb fixed hospital and neighborhood characteristics.
 \item County-by-year FE $\lambda_{ct}$ compare public and private hospitals in the same county-year and absorb county-wide reform and local shocks.
 \item County-clustered standard errors. Parallel trends are required for the public--private difference in net entry.
\end{itemize}
\end{frame}

\begin{frame}{Event-Study Specification and Timing}
\small
\[
 Y_{hctr}=\sum_{k\ne-1}\beta_{kr}
 Public_h\mathbb{1}\{t-T_c=k\}
 +\alpha_h+\lambda_{ct}+\varepsilon_{hctr}. \tag{5}
\]
\begin{itemize}\setlength\itemsep{0.6em}
 \item $T_c$ is the last pre-reform year, so $k=0$ is the first reformed year and $k=-1$ is the omitted reference.
 \item Endpoints bin four or more years before and after the reform.
 \item Joint lead tests use the estimable pre-reform indicators.
 \item Confidence intervals and joint tests use standard errors clustered by county.
\end{itemize}
\end{frame}
\begin{frame}{Net Pharmacy Entry at the Gate: Event Study}
\centering
\includegraphics[width=0.88\textwidth,height=0.70\textheight,keepaspectratio]{\paperresultsroot/figure_event_gate.pdf}
\par\footnotesize The vertical line separates the omitted pre-period from the adoption year.
The joint test of the pre-reform coefficients has $p=0.137$. The test does not establish parallel trends.
\end{frame}

\begin{frame}{Effects of the Zero-Markup Reform on Pharmacies by Distance}
\centering\scriptsize\setlength{\tabcolsep}{5.5pt}
\input{table_main_results_1028.tex}
\par\vspace{0.45em}
\raggedright\scriptsize
Notes: All regressions include hospital and county-by-year fixed effects. Standard errors clustered by county are in parentheses. $^{*}p<0.10$, $^{**}p<0.05$, $^{***}p<0.01$.
\end{frame}

\begin{frame}{Effects on Pharmacies by Distance Ring}
\centering
\includegraphics[width=0.88\textwidth,height=0.70\textheight,keepaspectratio]{\paperresultsroot/figure_ring_effects.pdf}
\par\footnotesize Circles: annual net entry. Diamonds: active pharmacies.
Bars show 95\% confidence intervals; effects are estimated directly on each ring.
\end{frame}

\begin{frame}{Gross Pharmacy Entry and Exit}
\centering
\small
\setlength{\tabcolsep}{8pt}
\begin{tabular}{lrrrr}
\toprule
& 200 m & 500 m & 800 m & 1,000 m \\
\midrule
Gross entry
& 0.043$^{***}$ & 0.068$^{**}$ & 0.006 & -0.052 \\
& (0.014) & (0.030) & (0.048) & (0.058) \\
Gross exit
& -0.027$^{**}$ & -0.043$^{*}$ & -0.001 & -0.018 \\
& (0.011) & (0.026) & (0.046) & (0.060) \\
\midrule
Outcome mean: entry & 0.227 & 1.019 & 2.215 & 3.163 \\
Outcome mean: exit & 0.122 & 0.567 & 1.243 & 1.785 \\
Observations & 93,265 & 93,265 & 93,265 & 93,265 \\
\bottomrule
\end{tabular}

\vspace{0.7em}
\footnotesize
The gate effect combines more pharmacy appearances with fewer disappearances. Neither margin changes at wider radii.
\par\vspace{0.35em}
\raggedright\scriptsize Notes: Entry is a \texttt{shop\_id} observed in year $t$ but not $t-1$; exit is defined conversely. These are platform appearance and disappearance measures. Models absorb hospital and county-by-year fixed effects; standard errors are clustered by county. $^{*}p<0.10$, $^{**}p<0.05$, $^{***}p<0.01$.
\end{frame}

\begin{frame}{Gross Pharmacy Entry: Event Studies}
\centering
\includegraphics[width=0.96\textwidth,height=0.69\textheight,keepaspectratio]{\allresultsroot/gross_event_entry_combined.pdf}
\par\footnotesize The omitted period is one year before the policy year. Each panel reports the joint test of the estimable pre-reform coefficients; bars show 95\% confidence intervals.
\end{frame}

\begin{frame}{Gross Pharmacy Exit: Event Studies}
\centering
\includegraphics[width=0.96\textwidth,height=0.69\textheight,keepaspectratio]{\allresultsroot/gross_event_exit_combined.pdf}
\par\footnotesize The omitted period is one year before the policy year. Each panel reports the joint test of the estimable pre-reform coefficients; bars show 95\% confidence intervals.
\end{frame}

\begin{frame}{Stacked DID: Effects by Distance Ring}
\centering\small\setlength{\tabcolsep}{7pt}
\input{table_stacked_rings_1028.tex}
\par\vspace{0.7em}
\raggedright\footnotesize
The mutually exclusive rings reproduce the baseline spatial pattern. Net entry rises at the gate and falls in the outer ring, leaving the 1 km total small and imprecise. The stock estimates show the same reallocation.
\par\vspace{0.35em}
\scriptsize Notes: Each stack pairs one public-hospital cohort with private hospitals in counties containing both groups. Models absorb stack-by-hospital and stack-by-county-by-year fixed effects. County-clustered standard errors are in parentheses. $^{*}p<0.10$, $^{**}p<0.05$, $^{***}p<0.01$.
\end{frame}

\begin{frame}{Robustness: Stacked DID with County-by-Year Fixed Effects}
\centering\small\setlength{\tabcolsep}{7pt}
\input{\paperresultsroot/table_stacked_did.tex}
\par\vspace{0.7em}
\raggedright\footnotesize
Each stack pairs one public-hospital treatment cohort with non-public hospitals in counties containing both groups. The model absorbs stack-by-hospital and stack-by-county-by-year fixed effects; standard errors are clustered by county. Net entry rises within 200 m and 500 m, then becomes small and imprecise at wider radii. Stock estimates are supporting evidence because their lead tests reject.
\end{frame}

\begin{frame}{Stacked DID: Net-Entry Event Studies}
\centering
\includegraphics[width=0.94\textwidth,height=0.68\textheight,keepaspectratio]{\paperresultsroot/figure_stacked_event_ne_200_500.pdf}
\par\vspace{0.35em}
\footnotesize The omitted period is one year before treatment. The vertical line marks the transition into treatment. Joint lead-test $p$-values are 0.102 within 200 m and 0.707 within 500 m; bars show 95\% confidence intervals.
\end{frame}

\begin{frame}{Robustness: Estimators and Treatment Definitions}
\centering\scriptsize\setlength{\tabcolsep}{5pt}
{\renewcommand{\arraystretch}{1.06}\input{table_robustness_1028.tex}}
\par\vspace{0.2em}
\raggedright\tiny
Notes: Annual net pharmacy entry within each radius. All models include hospital and county-by-year fixed effects. County-clustered standard errors are in parentheses. $^{*}p<0.10$, $^{**}p<0.05$, $^{***}p<0.01$.
\end{frame}

\begin{frame}{Placebo: Non-Pharmacy Wellness Establishments}
\centering
\small
\setlength{\tabcolsep}{7pt}
\begin{tabular}{lcccc}
\toprule
 & 0--200 m & 200--500 m & 500--1,000 m & 0--1,000 m \\
\midrule
Net entry
 & -0.026 & -0.038 & -0.005 & -0.069 \\
 & (0.017) & (0.033) & (0.056) & (0.073) \\
Stock
 & 0.020 & 0.039 & -0.158$^{**}$ & -0.099 \\
 & (0.013) & (0.030) & (0.062) & (0.083) \\
\bottomrule
\end{tabular}

\vspace{0.7em}
\footnotesize
\begin{itemize}\setlength\itemsep{0.45em}
 \item Net entry does not increase in any distance band; the total stock within 1 km also shows no significant increase.
 \item The negative stock estimate at 500--1,000 m also provides no evidence of expansion. Together, the results do not support a broad rise in health demand or ordinary commercial growth around hospitals.
\end{itemize}

\vspace{0.15em}
\raggedright\scriptsize
Notes: Outcomes count non-pharmacy wellness establishments. Net entry is the annual change in the count; stock is the annual level. All models absorb hospital and county-by-year fixed effects. Standard errors clustered by county are in parentheses. ** $p<0.05$.
\end{frame}

\sectiondivider{Concentration around Tertiary Hospitals}

\begin{frame}{Hospital-Tier Model and Formal Difference Tests}
\small
\[
 Y^{(r)}_{hct}=\beta_{1r}D_{hct}
 +\delta_{2r}D_{hct}\mathbb{1}\{G_h=2\}
 +\delta_{3r}D_{hct}\mathbb{1}\{G_h=3\}
 +\alpha_h+\lambda_{c\times t}+\varepsilon^{(r)}_{hct}.
\]
\begin{itemize}\setlength\itemsep{0.6em}
 \item $D_{hct}=Public_h\times Post_{ct}$; $G_h$ is the final hospital-tier classification.
 \item Total effects: Grade 1 $=\beta_{1r}$; Grade 2 $=\beta_{1r}+\delta_{2r}$; Grade 3 $=\beta_{1r}+\delta_{3r}$.
 \item The formal Grade 3--Grade 1 contrast is $\delta_{3r}$, computed from the pooled model with its full covariance matrix.
 \item Tier results use cumulative radii, matching Table 4 of the paper.
\end{itemize}
\end{frame}

\begin{frame}{Hospital Tier: Annual Net Pharmacy Entry}
\centering\scriptsize\setlength{\tabcolsep}{7pt}
\input{\paperresultsroot/table_tier_ne.tex}
\par\vspace{0.6em}
\raggedright\footnotesize Classified-tier sample; all public hospitals treated.
Hospital and county-by-year FE; county-clustered SE.
Tier includes administrative grades and classifications inferred from hospital scale.
\end{frame}

\begin{frame}{Hospital Tier: Active Pharmacy Stocks}
\centering\scriptsize\setlength{\tabcolsep}{7pt}
\input{\paperresultsroot/table_tier_st.tex}
\par\vspace{0.6em}
\raggedright\footnotesize Each radius is estimated in one pooled interaction model.
The last row formally tests the Grade 3--Grade 1 difference.
County-clustered SE in parentheses.
\end{frame}

\begin{frame}{Effects on Pharmacy Stocks by Hospital Tier and Radius}
\centering
\includegraphics[width=0.88\textwidth,height=0.70\textheight,keepaspectratio]{\paperresultsroot/figure_tier_st.pdf}
\par\footnotesize Cumulative radii; total effects from the pooled model.
Bars show 95\% confidence intervals. Horizontal offsets improve readability.
\end{frame}

\begin{frame}{Effects on Net Entry by Hospital Tier and Radius}
\centering
\includegraphics[width=0.88\textwidth,height=0.70\textheight,keepaspectratio]{\paperresultsroot/figure_tier_ne.pdf}
\par\footnotesize Grade 3 shows positive net entry; Grade 1 estimates turn negative at wider radii.
Bars show 95\% confidence intervals.
\end{frame}

\begin{frame}{Public versus Non-Public Grade 3 Hospitals}
\centering
\includegraphics[width=0.88\textwidth,height=0.70\textheight,keepaspectratio]{\paperresultsroot/figure_grade3_public_vs_nonpublic.pdf}
\par\footnotesize The coefficients are within-tier DID estimates: public Grade 3 hospitals are treated and non-public Grade 3 hospitals are the untreated comparison group. Bars show 95\% confidence intervals.
\end{frame}

\begin{frame}{Spatial Concentration and Access}
\small
\begin{itemize}\setlength\itemsep{0.8em}
 \item Within hospital neighborhoods: entry at the gate offsets losses in the outer ring, leaving the 1 km total imprecise.
 \item Across hospitals: supply gains concentrate around Grade 3 public hospitals, while wider-radius stocks fall around Grade 1 hospitals.
 \item Patients filling hospital prescriptions may gain convenient outside options; residents relying on neighborhood pharmacies may face longer trips.
 \item Location counts do not identify individual pharmacy relocations or measure patient welfare directly. These remain interpretation and future measurement tasks.
\end{itemize}
\end{frame}

\begin{frame}{Residential Proximity and Population-Weighted Access}
\footnotesize
\begin{columns}[T,onlytextwidth]
  \begin{column}{0.48\textwidth}
    {\color{AcademicBlue}\normalsize Pharmacy--community distance}\par\vspace{0.35em}
    \begin{itemize}\setlength\itemsep{0.4em}
      \item Residential communities come from annual Amap POI records for 2013--2018.
      \item For each pharmacy, we calculate the distance to its nearest residential community.
      \item Hospital-year means use cumulative radii (200, 500, and 1,000 m) and mutually exclusive rings (0--200, 200--500, and 500--1,000 m).
    \end{itemize}
  \end{column}
  \begin{column}{0.48\textwidth}
    {\color{AcademicBlue}\normalsize Spatial inequality in access}\par\vspace{0.35em}
    \begin{itemize}\setlength\itemsep{0.4em}
      \item WorldPop annual population rasters provide approximately 100 m residential demand cells.
      \item 2SFCA first divides each pharmacy by population in its 500 m catchment, then sums nearby ratios for each grid cell.
      \item Population-weighted Gini is calculated within 500 m, 2 km, and 3 km of each hospital. The 2 km radius is the main specification.
      \item Annual population weights are the main measure. Fixed 2013 weights hold population locations constant.
    \end{itemize}
  \end{column}
\end{columns}
\vfill
\tiny Sources: Amap POI (residential communities); WorldPop annual gridded population, 2013--2018. Hospital distance is measured from the AOI boundary when available and from the hospital point otherwise.
\end{frame}

\begin{frame}{Pharmacies Moved Closer to Communities in the 200--500 m Ring}
\centering
\footnotesize
\begin{tabular}{lrrrc}
\toprule
Hospital area & DID effect (m) & Standard error & $p$-value & Observations \\
\midrule
Stock: within 200 m       & -49.51 & 41.96 & 0.238 & 29,508 \\
Stock: within 500 m       & -55.08$^{*}$ & 30.68 & 0.073 & 50,036 \\
Stock: within 1,000 m     & -44.51 & 30.04 & 0.139 & 59,366 \\
\addlinespace
Ring: 0--200 m            & -49.51 & 41.96 & 0.238 & 29,508 \\
Ring: 200--500 m          & -58.51$^{**}$ & 26.20 & 0.026 & 46,205 \\
Ring: 500--1,000 m        & -18.33 & 18.36 & 0.318 & 54,688 \\
\bottomrule
\end{tabular}

\vspace{0.55em}
\begin{itemize}\setlength\itemsep{0.3em}
  \item Negative coefficients indicate that pharmacies around treated public hospitals became closer to their nearest residential community relative to pharmacies around private hospitals.
  \item The result is concentrated in the 200--500 m ring: the mean pharmacy--community distance fell by about 59 m.
\end{itemize}
\vspace{0.2em}
\raggedright\tiny Notes: Outcomes are annual changes in the hospital-year mean distance. Models absorb hospital and county-by-year fixed effects; standard errors are clustered by county. $^{*}p<0.10$, $^{**}p<0.05$.
\end{frame}

\begin{frame}{Relocating Pharmacies: Distances to Communities and Hospitals}
\centering
\scriptsize
\begingroup
\renewcommand{\arraystretch}{1.02}
\begin{tabular}{lrrr}
\toprule
Hospital area & $\Delta$ nearest community (m) & $\Delta$ nearest hospital (m) & Observations \\
\midrule
Stock: within 200 m
 & -28.42 & -98.00$^{*}$ & 360 \\
 & (29.91) & (55.72) & \\
Stock: within 500 m
 & -35.86 & -63.29 & 5,195 \\
 & (28.71) & (47.48) & \\
Stock: within 1,000 m
 & -5.19 & 30.71 & 19,541 \\
 & (16.88) & (31.19) & \\
\addlinespace
Ring: 0--200 m
 & -28.42 & -98.00$^{*}$ & 360 \\
 & (29.91) & (55.72) & \\
Ring: 200--500 m
 & -5.50 & -29.48 & 3,260 \\
 & (18.52) & (36.57) & \\
Ring: 500--1,000 m
 & -20.38 & 44.12 & 13,063 \\
 & (23.16) & (40.34) & \\
\bottomrule
\end{tabular}
\endgroup

\vspace{0.2em}
\begin{itemize}\setlength\itemsep{0.2em}
  \item At the hospital gate (0--200 m), relocating pharmacies moved about 98 m closer to their nearest hospital relative to the control group.
  \item Community-distance effects are negative and insignificant, so the results do not support relocation away from residential communities.
\end{itemize}
\vspace{0.05em}
\raggedright\tiny Notes: A relocating pharmacy is the same \texttt{shop\_id} moving at least 100 m between adjacent years. Outcomes are destination-minus-origin distance changes, evaluated against the same destination-year POI sets. DID models absorb hospital and county-by-year fixed effects; standard errors (in parentheses) are clustered by county. $^{*}p<0.10$.
\end{frame}

\begin{frame}{Population-Relative Pharmacy-Share Gini: No Robust Reform Effect}
\centering
\footnotesize
\begin{tabular}{lrrr}
\toprule
 & 500 m & 2 km (main) & 3 km \\
\midrule
Annual population weights
 & -0.00035 & 0.00002 & -0.00000 \\
 & (0.00035) & (0.00008) & (0.00006) \\
Fixed 2013 population weights
 & -0.00036 & 0.00002 & 0.00001 \\
 & (0.00035) & (0.00008) & (0.00006) \\
\addlinespace
Observations
 & 72,956 & 90,784 & 93,177 \\
Lead-test $p$-value (annual weights)
 & 0.534 & 0.656 & 0.117 \\
\bottomrule
\end{tabular}

\vspace{0.45em}
\begin{itemize}\setlength\itemsep{0.25em}
  \item The Gini compares each grid cell's pharmacy share with its population share; zero indicates proportional spatial matching and larger values indicate greater mismatch.
  \item The 2 km main estimate is economically tiny and statistically insignificant. Results are similar when population is fixed at its 2013 distribution.
\end{itemize}
\vspace{0.1em}
\raggedright\tiny Notes: WorldPop native population cells are retained, including populated cells with zero pharmacies. Cells are ranked by pharmacies per resident. Models absorb hospital and county-by-year fixed effects; standard errors (in parentheses) are clustered by county. No coefficient is statistically significant at the 10\% level.
\end{frame}

\begin{frame}{2SFCA Results: No Increase in Spatial Inequality}
\centering
\footnotesize
\begin{tabular}{lrrr}
\toprule
 & 500 m & 2 km (main) & 3 km \\
\midrule
Annual population weights
 & -0.0039$^{*}$ & -0.0017$^{*}$ & -0.0011 \\
 & (0.0023) & (0.0009) & (0.0007) \\
Fixed 2013 population weights
 & -0.0038$^{*}$ & -0.0015 & -0.0007 \\
 & (0.0023) & (0.0009) & (0.0007) \\
\addlinespace
Outcome mean (annual weights)
 & 0.3420 & 0.6016 & 0.6576 \\
Observations
 & 84,046 & 92,136 & 93,920 \\
Lead-test $p$-value
 & 0.624 & 0.823 & 0.373 \\
\bottomrule
\end{tabular}

\vspace{0.5em}
\begin{itemize}\setlength\itemsep{0.25em}
  \item All estimates are negative. The data provide no evidence that the reform increased inequality in population-adjusted pharmacy accessibility.
  \item In the 2 km main specification, Gini falls by 0.0017, about 0.3\% of its mean. The lead test provides no evidence of differential pre-trends.
\end{itemize}
\vspace{0.15em}
\raggedright\tiny Notes: 2SFCA accessibility is the sum of pharmacy-to-population ratios within a 500 m floating catchment, scaled per 10,000 residents. The outcome is its population-weighted Gini within each hospital radius. Models absorb hospital and county-by-year fixed effects; standard errors are clustered by county. $^{*}p<0.10$.
\end{frame}

\sectiondivider{Heterogeneity}

\begin{frame}{County Characteristics and the Net-Entry Effect}
\centering
\scriptsize
\setlength{\tabcolsep}{6.5pt}
\begin{tabular}{lrrrr}
\toprule
DID $\times$ standardized characteristic & 200 m & 500 m & 800 m & 1,000 m \\
\midrule
Population age 65+ share
& 0.027 & 0.092 & 0.119 & 0.276$^{**}$ \\
& (0.027) & (0.059) & (0.094) & (0.115) \\
Urban population share
& 0.045 & 0.094$^{*}$ & 0.097 & 0.162 \\
& (0.028) & (0.055) & (0.091) & (0.118) \\
Average years of education
& 0.065$^{**}$ & 0.128$^{**}$ & 0.160 & 0.226$^{*}$ \\
& (0.029) & (0.065) & (0.104) & (0.131) \\
Log population density
& 0.062$^{**}$ & 0.138$^{**}$ & 0.179$^{*}$ & 0.234$^{**}$ \\
& (0.027) & (0.057) & (0.092) & (0.118) \\
\midrule
Observations & 67,546 & 67,546 & 67,546 & 67,546 \\
\bottomrule
\end{tabular}

\vspace{0.55em}
\footnotesize The response is stronger in denser and more educated counties. Urbanization is weaker, while the elderly-share interaction appears only at 1,000 m.
\par\vspace{0.25em}
\raggedright\tiny Notes: Each row comes from a separate model. Characteristics are fixed at their 2010 census values and standardized in the estimation sample. Models include the DID main effect, hospital FE, and county-by-year FE; standard errors are clustered by county. The correlated characteristics and multiple tests make these exploratory results. $^{*}p<0.10$, $^{**}p<0.05$.
\end{frame}

\begin{frame}{General and Specialist Hospitals}
\centering
\small
\setlength{\tabcolsep}{7pt}
\begin{tabular}{lrrrr}
\toprule
Net-entry effect & 200 m & 500 m & 800 m & 1,000 m \\
\midrule
General hospitals
& 0.227$^{***}$ & 0.311$^{***}$ & 0.206$^{*}$ & 0.162 \\
& (0.042) & (0.080) & (0.120) & (0.158) \\
Specialist hospitals
& 0.091$^{*}$ & 0.010 & -0.155 & -0.351 \\
& (0.051) & (0.104) & (0.178) & (0.229) \\
Specialist minus general
& -0.136$^{**}$ & -0.301$^{**}$ & -0.360$^{*}$ & -0.513$^{**}$ \\
& (0.066) & (0.127) & (0.199) & (0.254) \\
\midrule
Observations & 55,570 & 55,570 & 55,570 & 55,570 \\
\bottomrule
\end{tabular}

\vspace{0.7em}
\footnotesize Pharmacy growth is concentrated around general hospitals. Specialist-hospital effects are small or negative, and the formal specialist-minus-general contrasts are negative at every radius.
\par\vspace{0.25em}
\raggedright\scriptsize Notes: The pooled interaction model restricts the sample to hospitals classified as general or specialist. It includes hospital and county-by-year fixed effects; standard errors are clustered by county. $^{*}p<0.10$, $^{**}p<0.05$, $^{***}p<0.01$.
\end{frame}

\sectiondivider{Conclusion}
\begin{frame}{Main Findings}
\small
\begin{enumerate}\setlength\itemsep{0.55em}
  \item Hospital net entry rises in the 0--200m gate band and falls in the 500--1,000m outer ring; the 1km total is small and imprecise. Gains concentrate around Grade 3 hospitals.
  \item At the gate, the reform increases gross pharmacy entry and reduces gross exit. Neither margin changes at wider radii.
  \item The pharmacy response is stronger in denser and more educated counties, but correlated characteristics and multiple tests make the interaction results exploratory.
  \item Net entry increases around general hospitals and shows no clear positive response around specialist hospitals.
  \item Population-adjusted access measures provide no evidence that spatial inequality increased after the reform.
\end{enumerate}
\end{frame}

\end{document}
